\section{Preliminaries}
The \emph{Boolean semiring} is the semiring on $\{0,1\}$ with OR as addition and AND as multiplication. For Boolean matrices $A$ and $B$, $A \vee B$ is the componentwise OR of $A$ and $B$, $A \wedge B$ is the componentwise AND, and $A \star B$ is the (Boolean) matrix product over the Boolean semiring. When it is clear from the context, sometimes we omit the $\star$ and write $A B$ for the product.

Since the running times of our algorithms involve polylogarithmic terms, we must make the computational model precise. Unless otherwise specified, we assume a standard word RAM with wordsize $w$. That is, accessing a memory location takes $O(1)$ time, and we can perform simple operations (such as addition, componentwise AND and XOR, but not multiplication) on $w$-bit numbers in $O(1)$ time.
Typically, speedups in combinatorial algorithms come from either exploiting some combinatorial substructure, by preprocessing and doing table lookups, or by some ``word tricks'' which utilize the bit-level parallelism of the machine model.
In our results, we explicitly state the dependence of the word size, denoted by $w$. The reader may assume $w = \Theta(\log n)$ for convenience. In fact all algorithms in this paper can be implemented on a pointer machine under this constraint.

We now describe some of the tools we need.
